% Vietnamese translation for OI Public Library
% Source-File: IOI中国国家候选队论文/2005/黄刚/黄刚.ppt
\documentclass[11pt]{article}
\usepackage{oipl-vn}

\newcommand{\Oh}{\mathcal{O}}

\title{Bản trình chiếu: Dynamic Rankings}
\author{Huang Gang}
\date{2005}

\begin{document}
\maketitle

\origtitle{Dynamic Rankings: bản trình chiếu}
\sourcepath{IOI China candidate team papers / 2005 / Huang Gang / Huang Gang.ppt}

\section{Mạch trình bày}

Tệp gốc chỉ có PowerPoint, nhưng phần chữ đọc được từ slide khá rõ. Chủ đề là
\textit{Dynamic Rankings}; giới hạn xuất hiện trên slide gồm
\[
  1\le N\le 50000,\qquad 1\le A[i]\le \mathrm{MaxLongint},
\]
cùng các thao tác \(C(i,j)\), \(Q(i,j,k)\), \(T(i,j,w)\). Các mốc độ phức
tạp gồm \(\Oh(\log N)\), \(\Oh(N\log N)\), \(\Oh(\log^2 N)\) và tổng
\(\Oh(M\log\mathrm{MaxLongint}\log^2 N)\). Bản ghi chú này dịch lại mạch
slide: từ truy vấn xếp hạng tĩnh, thêm cập nhật, rồi quy truy vấn hạng về
truy vấn đếm.

\section{Bài toán}

Ta có dãy
\[
  A[1],A[2],\ldots,A[N].
\]
Truy vấn
\[
  Q(i,j,k)
\]
yêu cầu tìm phần tử nhỏ thứ \(k\) trong đoạn
\[
  A[i],A[i+1],\ldots,A[j].
\]
Thao tác
\[
  C(i,j)
\]
cập nhật \(A[i]\) thành \(j\). Slide có ví dụ nhỏ
\[
  N=6,\qquad A=[3,1,6,5,2,4],
\]
để minh họa rằng truy vấn không hỏi thứ tự toàn cục của cả dãy, mà hỏi thứ tự
của một đoạn con thay đổi theo thời gian.

Nếu không có cập nhật, bài toán có thể xử lý offline hoặc xây cấu trúc dữ liệu
tĩnh. Nếu chỉ hỏi hạng trên toàn bộ dãy, một cây cân bằng hoặc cây Fenwick
theo giá trị cũng đủ. Điểm khó của \textit{Dynamic Rankings} là hai chiều động
cùng xuất hiện: đoạn truy vấn thay đổi theo chỉ số, còn giá trị tại một vị trí
cũng thay đổi sau các lệnh \(C\).

\section{Mô hình dữ liệu}

Slide dẫn từ truy vấn toàn cục \(Q(k)\), rồi đến truy vấn đoạn \(Q(i,j)\), và
cuối cùng tới \(Q(i,j,k)\). Cách nhìn này rất quan trọng: trước khi tìm phần
tử hạng \(k\), ta cần biết cách trả lời câu hỏi đơn giản hơn:
\[
  T(i,j,w)=\#\{x\mid i\le x\le j,\ A[x]\le w\}.
\]
Nói cách khác, \(T(i,j,w)\) đếm có bao nhiêu phần tử trong đoạn \([i,j]\)
không vượt quá ngưỡng \(w\). Ví dụ slide ghi dãy
\[
  A=[1,4,2,1]
\]
và truy vấn \(T(2,4,3)\). Đoạn cần xét là \(4,2,1\), nên có hai phần tử không
vượt quá \(3\).

Khi đã có \(T\), truy vấn hạng được chuyển thành tìm ngưỡng nhỏ nhất \(w\) sao
cho
\[
  T(i,j,w)\ge k.
\]
Vì \(T(i,j,w)\) đơn điệu theo \(w\), có thể nhị phân trên miền giá trị. Đây là
ý tưởng trung tâm của slide: không trực tiếp duy trì phần tử hạng \(k\), mà
duy trì một hàm đếm đủ mạnh để khôi phục hạng bằng tìm kiếm.

\section{Cấu trúc hai tầng}

Một hiện thực chuẩn là cây đoạn hoặc cây Fenwick theo chỉ số, trong mỗi nút
lưu một cấu trúc có thứ tự theo giá trị. Với một tiền tố \(A[1..p]\), Fenwick
tree biểu diễn nó bằng \(\Oh(\log N)\) nút. Vì thế đoạn \([i,j]\) được trả lời
bằng hiệu của hai tiền tố \(j\) và \(i-1\). Trong từng nút, ta cần biết có bao
nhiêu giá trị không vượt quá \(w\); việc này có thể làm bằng cây cân bằng,
mảng đã sắp xếp nếu tĩnh, hoặc cấu trúc động tương ứng nếu có cập nhật.

Với cập nhật \(C(i,j)\), giá trị cũ của \(A[i]\) bị xóa khỏi mọi nút Fenwick
chứa vị trí \(i\), và giá trị mới \(j\) được chèn vào các nút ấy. Có
\(\Oh(\log N)\) nút bị ảnh hưởng. Nếu mỗi thao tác chèn, xóa, đếm trong cấu
trúc con tốn \(\Oh(\log N)\), một lần cập nhật hoặc một lần tính \(T(i,j,w)\)
tốn \(\Oh(\log^2 N)\). Đây là mốc độ phức tạp xuất hiện trong PowerPoint.

Slide cũng ghi mốc xây dựng \(\Oh(N\log N)\). Điều này phù hợp với cách đưa
mỗi phần tử ban đầu vào các nút Fenwick/cây đoạn tương ứng. Nếu cần nén giá
trị, ta có thể thu thập trước mọi giá trị ban đầu và mọi giá trị xuất hiện
trong cập nhật, rồi đánh số lại; khi đó nhị phân theo giá trị chạy trên tập
giá trị rời rạc thay vì toàn bộ \(\mathrm{MaxLongint}\).

\section{Từ T tới Q}

Để trả lời \(Q(i,j,k)\), ta nhị phân \(w\). Với mỗi \(w\), gọi
\[
  c=T(i,j,w).
\]
Nếu \(c\ge k\), phần tử hạng \(k\) không lớn hơn \(w\), nên thu hẹp nửa trái.
Nếu \(c<k\), cần tìm giá trị lớn hơn. Sau khi vòng nhị phân kết thúc, giá trị
nhỏ nhất thỏa \(T(i,j,w)\ge k\) chính là đáp án.

Nếu nhị phân trên toàn miền số nguyên, mỗi truy vấn tốn
\[
  \Oh(\log\mathrm{MaxLongint}\cdot\log^2 N).
\]
Với \(M\) lệnh, tổng thời gian khớp với mốc trên slide:
\[
  \Oh(M\log\mathrm{MaxLongint}\log^2 N),
\]
chưa tính bước xây dựng ban đầu. Nếu nén giá trị, thừa số
\(\log\mathrm{MaxLongint}\) được thay bằng \(\log V\), trong đó \(V\) là số
giá trị khác nhau có thể xuất hiện.

\section{Nhận xét}

Điểm đáng học của bản trình chiếu không chỉ là một cấu trúc dữ liệu cụ thể.
Nó là cách tách bài toán thống kê thứ tự thành hai lớp: lớp ngoài xử lý đoạn
theo chỉ số, lớp trong xử lý thứ tự theo giá trị. Truy vấn hạng \(k\) vốn có
vẻ khó cập nhật trực tiếp được thay bằng truy vấn đếm \(T(i,j,w)\), một câu
hỏi có tính cộng trừ theo tiền tố và đơn điệu theo \(w\).

Khi gặp biến thể khác của bài toán, ta nên kiểm tra ba điểm: thao tác cập nhật
có thể sửa bao nhiêu nút lớp ngoài; cấu trúc lớp trong có hỗ trợ đếm theo
ngưỡng không; và miền giá trị có thể nén được không. Ba câu hỏi này thường
quyết định lời giải đạt \(\Oh(\log^2 N)\) hay bị đội lên bởi các thao tác đếm
hoặc nhị phân không cần thiết.

\end{document}
